% CLAUDE REPORT::BRIDGE-DICHOTOMY-0001::2026-11-10T12:48-08:00
\documentclass[11pt]{article}
\usepackage{amsmath,amssymb,amsthm}
\usepackage[margin=1in]{geometry}
\usepackage{url}
\newtheorem{theorem}{Theorem}
\newtheorem{corollary}[theorem]{Corollary}
\newtheorem{lemma}[theorem]{Lemma}
\theoremstyle{definition}
\newtheorem{definition}[theorem]{Definition}
\newtheorem{remark}[theorem]{Remark}
\newcommand{\Wbar}{\overline W}
\newcommand{\TAUT}{\mathrm{TAUT}}
\title{The fixed-calculus bridge dichotomy for $\Wbar_R$:\\
why ``$\mathrm{SAT}\in\mathrm P\Rightarrow\Wbar_R(N)\le\mathrm{poly}(N)$'' is never weaker than the target when $R$ is weak, and what makes it true when $R$ is strong}
\author{Claude scheduled research note (supporting reference; not an LSPV role report)}
\date{10 November 2026, 12:48 PST}
\begin{document}
\maketitle

\paragraph{Status line.} Theorems~\ref{thm:collapse}, \ref{thm:monotone}, \ref{thm:optimal} are \textbf{PROVED/EXACT} (short proofs below, relying on cited \textbf{KNOWN THEOREMS}). Corollary~\ref{cor:mix} is \textbf{CONDITIONAL} on the corpus' still uncertified DEF checker semantics of $R_{\mathrm{mix341}}$. No superpolynomial lower bound on $\Wbar_{R_{\mathrm{mix341}}}$, no bridge B3 for $R_{\mathrm{mix341}}$, and no separation is claimed. \textbf{P versus NP remains OPEN.}

\section{Setting}
This note addresses the shared worst-case directive
\[
\Wbar_R(N)=\max_{F\in\mathrm{UNSAT},\ |\langle F\rangle|\le N}\ \min_{\pi:F\vdash_R\bot}\operatorname{cost}_R(\pi),
\qquad
\textbf{Bridge}(R):\ \mathrm{SAT}\in\mathrm P\Longrightarrow \Wbar_R(N)\le N^{O(1)} .
\]
The multisign DAG-hypergraph reduction system of the directive (sign-merging of hyperedges, with the minimum taken over all valid reduction sequences and with the optimal sequence granted for free) is, once its rules are fixed, a propositional proof system $R$ in the sense of Cook--Reckhow: a polynomial-time checker $V_R(F,\pi)$ that accepts only when $F$ is unsatisfiable, with $\operatorname{cost}_R(\pi)$ polynomially related to $|\pi|$. Throughout, $R$ is assumed to be such a system (sound, complete, polynomially checkable, cost polynomially equivalent to proof length). All earlier corpus calculi (ordinary Resolution, $R_{\mathrm{mix341}}$, the identity system $Q_A$, $R_T^{\mathrm{len}}$) are of this kind.

Granting the optimal derivation for free is correct for the \emph{lower-bound} direction, as the directive says: $\min_\pi$ is a lower bound on the cost of any procedure that \emph{outputs} an $R$-derivation. The point of this note is about the \emph{other} arrow, Bridge$(R)$, which the directive calls the weakest sufficient bridge.

\section{The collapse theorem}
\begin{theorem}[Bridge collapse; PROVED/EXACT]\label{thm:collapse}
Let $R$ be a fixed proof system for which $\Wbar_R(N)\ge N^{\omega(1)}$ is a theorem. Then
\[
\textbf{Bridge}(R)\iff \mathrm{P}\neq\mathrm{NP}.
\]
\end{theorem}
\begin{proof}
($\Leftarrow$) If $\mathrm P\neq\mathrm{NP}$ the premise of Bridge$(R)$ is false, so the implication holds vacuously.
($\Rightarrow$) Bridge$(R)$ and the theorem $\neg(\Wbar_R\le\mathrm{poly})$ give $\mathrm{SAT}\notin\mathrm P$ by contraposition.
\end{proof}

\begin{remark}
The theorem is trivial, but its consequence is not. Haken (1985) and its successors (Urquhart 1987 Tseitin; Chv\'atal--Szemer\'edi 1988 random 3-CNF; Ben-Sasson--Wigderson 2001 width) make Theorem~\ref{thm:collapse} applicable to ordinary Resolution, and later results to Res$(k)$, cutting planes, polynomial calculus, bounded-depth Frege and, per the September 2026 Braun preprint cited in the corpus, unrestricted DAG Res$(\oplus)$. For every such $R$, proving Bridge$(R)$ \emph{is} proving $\mathrm P\ne\mathrm{NP}$. In particular, it cannot be ``weaker'' than the original goal. Every transfer theorem for such an $R$ must, as a by-product, give an unconditional superpolynomial lower bound for every SAT algorithm. So it is subject to the full relativization, natural-proofs and algebrization barriers, which the bridge formulation had seemed to avoid. This also explains structurally the corpus' repeated finding (LSPV-1/367 recognizer bypass, LSPV-5/289, LSPV-1/339 clock-gate) that a polynomial-time SAT algorithm can ignore a weak calculus. Under $\mathrm P\neq\mathrm{NP}$ nothing forbids that, and under $\mathrm P=\mathrm{NP}$ a fast algorithm on a hard-for-$R$ family is exactly what the contrapositive excludes.
\end{remark}

\begin{theorem}[Monotonicity; PROVED/EXACT]\label{thm:monotone}
If $R'$ polynomially simulates $R$ in size ($\Wbar_{R'}(N)\le \Wbar_R(N)^{O(1)}+N^{O(1)}$), then Bridge$(R)\Rightarrow$ Bridge$(R')$, and any superpolynomial lower bound for $R'$ transfers to $R$.
\end{theorem}
\begin{proof}Immediate by composing the polynomial bounds.\end{proof}

Theorems~\ref{thm:collapse}--\ref{thm:monotone} give a \textbf{dichotomy}. Moving $R$ up the simulation order makes Bridge$(R)$ easier and the lower bound harder. Moving it down does the opposite. As soon as a superpolynomial lower bound is proved, the bridge is \emph{literally} $\mathrm P\ne\mathrm{NP}$. A route that is ``bridge plus lower bound'' only gains something for an $R$ that \emph{currently has no lower bound}. Then the two halves are independent problems, and the lower-bound half is at least Extended-Frege hard whenever $R$ simulates ER (next section).

\section{Where Bridge(R) is actually provable: optimal systems}
\begin{definition}A proof system $R$ is \emph{optimal} if for every proof system $Q$ there is a polynomial $p$ with $\min_{\pi:F\vdash_R\bot}|\pi|\le p(\min_{\rho:F\vdash_Q\bot}|\rho|+|F|)$ for all UNSAT $F$. It is \emph{p-optimal} if moreover the translation is polynomial-time computable.\end{definition}

\begin{theorem}[Optimal systems satisfy the bridge; PROVED/EXACT from KNOWN THEOREMS]\label{thm:optimal}
\begin{enumerate}
\item[(a)] If $R$ is optimal then Bridge$(R)$ holds.
\item[(b)] If $\mathrm{SAT}\in\mathrm P$ then a p-optimal proof system exists, and every optimal system $R$ then has $\Wbar_R(N)\le N^{O(1)}$.
\item[(c)] Nevertheless, (b) does not yield Bridge$(R)$ for any \emph{preselected} $R$. The system produced in (b) is built from the hypothesised E-algorithm for an NE-complete set and is not known in advance.
\end{enumerate}
\end{theorem}
\begin{proof}
(a) Assume $\mathrm{SAT}\in\mathrm P$ via a correct polynomial-time decider $A$. The identity system $Q_A$ (proof of $F$ is $F$ itself, accepted iff $A$ rejects $F$; LSPV-5/409) is a sound, complete proof system with $\Wbar_{Q_A}(N)\le N$. Optimality of $R$ gives $\Wbar_R(N)\le p(2N)$.

(b) $\mathrm P=\mathrm{NP}\Rightarrow\mathrm E=\mathrm{NE}$ by padding. For $L\in\mathrm{NE}$ decided in nondeterministic time $2^{cn}$, the padded set $\{x\#1^{2^{|x|}}:x\in L\}$ lies in NP, hence in P, hence $L\in\mathrm{DTIME}(2^{O(n)})$. Krajíček--Pudl\'ak (1989, J. Symb. Logic 54) prove that $\mathrm E=\mathrm{NE}$ implies a p-optimal proof system for TAUT exists. Since $\mathrm{SAT}\in\mathrm P$ also implies $\mathrm{NP}=\mathrm{coNP}$, there is a polynomially bounded system ($Q_A$), and optimality transfers the polynomial bound as in (a).

(c) The Krajíček--Pudl\'ak construction enumerates proof systems whose soundness is decided by the E-time algorithm for an NE-complete problem. That algorithm exists only under the hypothesis. No fixed, explicitly given $R$ is currently known to be optimal. Existence of optimal systems is the open Krajíček--Pudl\'ak question, still open as of the 2024--2026 literature checked this run (Köbler--Messner--Torán 2003 weaken the hypothesis to $\mathrm{NEE}=\mathrm{coNEE}$ / $\mathrm{EE}=\mathrm{NEE}$; Pudl\'ak's finite-domain incompleteness conjectures (2017) predict that none exist).
\end{proof}

\begin{remark}[relation to corpus]
LSPV-5/409 proved Bridge$(R)\iff$ ``$R$ size-simulates $Q_A$'' for the hypothesised $A$, and warned that this is circular if it is simply assumed. LSPV-1/337 gave a conditional route via $T$-provable finite soundness ($R_T^{\mathrm{len}}$). LSPV-1/415 proved a $\Pi^0_1$-completeness obstruction to blanket effective checking. Theorem~\ref{thm:optimal} places all three inside one known frame. A preselected system satisfies the bridge for every possible $A$ at once exactly when it simulates every $Q_A$. Optimality is the natural, \emph{$A$-independent} sufficient condition, and it is an open, well-studied problem. The new content relative to the corpus is (i) Theorem~\ref{thm:collapse} stated explicitly for $\Wbar_R$ and (ii) the $\mathrm P=\mathrm{NP}\Rightarrow\mathrm E=\mathrm{NE}\Rightarrow$ p-optimal step. A repository search (\texttt{E=NE}, \texttt{NE}, \texttt{p-optimal}, \texttt{optimal proof system}) found only an LSPV-1/391 note that GitHub searches for p-optimal systems found no applicable compiler. No prior corpus item states (i) or (ii).
\end{remark}

\section{Consequence for the shared calculus}
\begin{corollary}[CONDITIONAL]\label{cor:mix}
Assume the charged two-input DEF rule of $R_{\mathrm{mix341}}$, together with RES, admits the standard extension-variable clause references, as interpreted in LSPV-1/341 (``BPHP to unary ER escape''), so that $R_{\mathrm{mix341}}$ polynomially simulates Extended Resolution, which is p-equivalent to Extended Frege. Then:
\begin{enumerate}
\item any proof of $\Wbar_{R_{\mathrm{mix341}}}(N)\ge N^{\omega(1)}$ is a superpolynomial Extended Frege lower bound, a central open problem with no known approach (no 2026 Frege/EF lower bound was found in this run's literature check);
\item if such a proof were obtained, Theorem~\ref{thm:collapse} would make the remaining B3 bridge for $R_{\mathrm{mix341}}$ equivalent to $\mathrm P\neq\mathrm{NP}$.
\end{enumerate}
If instead the DEF semantics are \emph{weaker} than ER, and $R_{\mathrm{mix341}}$ falls below a system with a proved exponential lower bound, then item~2 applies immediately and B3 is equivalent to $\mathrm P\neq\mathrm{NP}$ already.
\end{corollary}
Either way, the five-role program cannot get a proof of P$\neq$NP from ``hard family + fixed-calculus bridge'' that is easier than one of: (a) an EF-strength lower bound together with an independent bridge, or (b) P$\neq$NP itself. This is a calibration, not a falsification. It retires the label ``weakest'' for B3 when the calculus has a known lower bound, and it should be listed in the barrier ledger.

\section{Hostile controls and transfer classification}
\begin{itemize}
\item \textbf{Strategy-free minimum (directive item 3).} Free choice of $\pi$ does not affect Theorem~\ref{thm:collapse}. The theorem concerns the bridge arrow, not the minimum.
\item \textbf{Alternative reasoning (directive item 4).} Theorem~\ref{thm:monotone} is the exact form of ``hard instances must survive alternative reasoning''. Surviving every system is equivalent to $\mathrm{NP}\neq\mathrm{coNP}$ (Cook--Reckhow 1979), which is strictly stronger than what is needed. Surviving the single system $Q_A$ for every $A$ is what Bridge needs.
\item \textbf{Relativization.} The collapse argument is purely logical and relativizes. It therefore offers no new barrier evasion. It only shows that the bridge carries the whole difficulty.
\item \textbf{PQ ledger.} Item 5 (LWE worst-to-average-case architecture) is the closest analogue. Bridge$(R)$ asks that an arbitrary solver's behaviour be ``externally extracted'' as $R$-proofs. That is exactly simulation of $Q_A$. Classification: ANALOGY/CANDIDATE, with the missing theorem being an $A$-uniform $Q_A\to R$ compiler. Items 1--4 and 6--7: FAILED TRANSFER (no object correspondence at the bridge arrow).
\item \textbf{Cross-domain.} Optimization duality (weak duality certificates versus algorithmic optimality) is a useful analogy. A dual certificate bounds every primal algorithm only within the same LP. That is the same ``fixed system'' gap, and it is a misleading analogy if it is read as transferring across models.
\end{itemize}

\section{Corpus reviewed; duplication, contradiction, staleness}
\textbf{CORPUS REVIEWED:} \texttt{AGENTS.md} (including the 2026-11-09 knot/link handoff pointer and the ALL-ARC allocation); root \texttt{main.tex} regions on $R_{\mathrm{mix341}}$ (lines near LSPV-1/337, 339, 347, 367, 415, 539 entries); LSPV-1/341 ER escape; LSPV-1/391; the LSPV-5/465 prewrite ledger entry [P2-L5-409] ($Q_A$ identity boundary); the claude/ folder (TPORT, ONETAPE, world ledger); git log through commit \texttt{fa94368} (LSPV-4/524).
\textbf{DUPLICATION / CONTRADICTION AUDIT:} Consistent with LSPV-5/409, LSPV-1/337 and LSPV-1/415, and complementary to them. No contradiction found. Theorem~\ref{thm:collapse} sharpens the review note's existing remark that ``this bridge is classically equivalent to P$\neq$NP'' (\texttt{research/review/reports/proof\_map/2026-10-10\_INITIAL\_0001.md}) from the wrapper setting to $\Wbar_R$.
\textbf{WHAT CHANGED RELATIVE TO ROOT MAIN:} Nothing in \texttt{main.tex} was edited. This scheduled Claude run has no ownership region there, so the note is a supporting reference only.
\textbf{ADOPTED / REJECTED / RETIRED:} Adopted $Q_A$ (LSPV-5/409). Proposed retirement of the descriptor ``weakest sufficient'' for B3 whenever $R$ has a proved lower bound.
\textbf{STALE OR OVERSTATED CLAIMS FOUND:} None in the files read. The directive phrase ``weakest sufficient bridge'' is accurate only relative to stronger simulation theorems. It is not weaker than P$\neq$NP for weak calculi.

\section*{Literature checked this run}
Krajíček--Pudl\'ak, J. Symbolic Logic 54 (1989); Köbler--Messner--Torán, Inf. Comput. 184 (2003); Messner (E-hard sets lack p-optimal systems); arXiv:2408.07408 (2024, optimal proof systems for sets outside NQP); arXiv:2602.02294 (2026, no set in PSPACE$\setminus$NP has optimal proof systems, as summarized by the search; not independently verified); Egidy--Glaßer--Herold MFCS 2023 (Boolean hierarchy); Davis--Robere ECCC TR26-055 (2026, bounded-depth Frege proves its own lower bounds); Khaniki--Pich--Sokolov CCC 2026 (prover-adversary games, EF anticheckers); Alekseev, Theory of Computing 22 (2026, PC with extension); arXiv:2506.16956 (proof analysis for EF is NP-complete). None of these gives an explicit optimal proof system or a Frege/EF lower bound. Theorem~\ref{thm:optimal}(c) and Corollary~\ref{cor:mix}(1) remain OPEN obstacles.

\section*{Next targets}
\begin{enumerate}
\item OPEN: decide whether $R_{\mathrm{mix341}}$ DEF is ER-strength (LSPV-2 checker obligation). This fixes which side of the dichotomy the program is on.
\item OPEN, CANDIDATE: an $A$-uniform compiler $Q_A\to R$ for algorithms $A$ in a \emph{restricted} class (e.g. polynomial-time algorithms whose correctness is provable in $S^1_2$). This is the Cook--Krajíček $R_T$ direction and is the only known non-vacuous bridge shape.
\end{enumerate}

\bigskip\noindent P versus NP remains OPEN.
% END_OF_VERIFIED_REPORT::CLAUDE::BRIDGE-DICHOTOMY-0001
\end{document}
